

\section{CutMix}
\label{sec:cutmix}
\noindent
We describe the CutMix algorithm in detail. 

\subsection{Algorithm}
\label{sec:cutmix_algorithm}
Let $x \in \mathbb{R}^{W \times H \times C}$ and $y$ denote a training image and its label, respectively.
The goal of CutMix is to generate a new training sample $(\Tilde{x},\Tilde{y})$ by combining two training samples $(x_{A}, y_{A})$ and $(x_{B}, y_{B})$. 
{The generated} training sample $(\Tilde{x},\Tilde{y})$ is used to train the model {with} its original loss function.
{We define the combining operation as}
\begin{equation}
\begin{split}
    \Tilde{x} & =  \mathbf{M} \odot x_{A} + (\mathbf{1}- \mathbf{M}) \odot x_{B} \\
    \Tilde{y} & =  \lambda y_A + (1-\lambda) y_B, \\
\end{split}
\label{eq:cutmix}
\end{equation}
where $\mathbf{M} \in \{0,1\}^{W \times H}$ denotes a binary mask indicating where to drop out and fill in from two images, $\mathbf{1}$ is a binary mask filled with ones, and $\odot$ is element-wise multiplication.
Like Mixup~\cite{zhang2017mixup}, the combination ratio $\lambda$ between two data points is sampled from the beta distribution $\text{Beta}(\alpha,\alpha)$. 
In our all experiments, we set $\alpha$ to $1$, that is $\lambda$ is sampled from the uniform distribution $(0,1)$. 
Note that the major difference is that CutMix replaces an image region with a patch from another training image and generates {more locally natural image than Mixup does}.

To sample the binary mask $\mathbf{M}$, we first sample the bounding box coordinates $\mathbf{B}=(r_x,r_y,r_w,r_h)$ indicating the cropping regions on $x_A$ and $x_B$. 
The region $\mathbf{B}$ in $x_A$ is removed and filled in with the patch cropped from $\mathbf{B}$ of $x_B$. 


In our experiments, we sample rectangular masks $\mathbf{M}$ whose aspect ratio is proportional to the original image. The box coordinates are uniformly sampled according to:  
\begin{equation}
\begin{split}
    r_{x} \sim \text{Unif}~ \left(0, W\right), &~~~r_{w} = W \sqrt{1-\lambda},  \\
    r_{y} \sim \text{Unif}~ \left(0, H\right), &~~~r_{h} = H \sqrt{1-\lambda}  \\
\end{split}
\label{eq:sampling}
\end{equation}
making the cropped area ratio $\frac{r_w r_h}{WH} = 1-\lambda$.
With the cropping region, the binary mask $\mathbf{M}$ $\in \{0,1\}^{W \times H}$ is decided by filling with $0$ within the bounding box $\mathbf{B}$, otherwise $1$.

In each training iteration, a CutMix-ed sample $(\Tilde{x}, \Tilde{y})$ is generated by combining randomly selected two training samples in a mini-batch according to Equation~(\ref{eq:cutmix}).
Code-level details are presented in Appendix~\ref{appendix:algorithm}. 
CutMix is simple and incurs a negligible computational overhead {as} existing data augmentation techniques used in~\cite{InceptionResnet,densenet}; {we can efficiently utilize it to train any network architecture.}

\subsection{Discussion}
\label{sec:cutmix_discussion}


\noindent\textbf{What does model learn with CutMix? }
We have motivated CutMix such that full object extents are considered as cues for classification, the motivation shared by Cutout, while ensuring two objects are recognized from partial views in a single image to increase training efficiency.
To verify that CutMix is indeed learning to recognize two objects from their respective partial views, we visually compare the activation maps for CutMix against Cutout~\cite{devries2017cutout} and Mixup~\cite{zhang2017mixup}. 
Figure~\ref{fig:class_cam} shows example augmentation inputs as well as corresponding class activation maps (CAM)~\cite{zhou2016CAM} for two classes present, Saint Bernard and Miniature Poodle. 
We {use} vanilla ResNet-50 model\footnote{We use ImageNet-pretrained ResNet-50 provided by PyTorch~\cite{paszke2017automatic}.} for obtaining the CAMs to clearly see the effect of augmentation method only.


\begin{figure}
    \centering
    \includegraphics[width=0.77\linewidth]{cam8.pdf}
    \caption{
    Class activation mapping (CAM)~\cite{zhou2016CAM} visualizations on `Saint Bernard' and `Miniature Poodle' samples using various augmentation techniques.
    From top to bottom rows, we show the original images, input augmented image, CAM for class `Saint Bernard', and CAM for class `Miniature Poodle', respectively. 
    {Note that CutMix can take advantage of the mixed region on image, but Cutout cannot.}
    }
    \label{fig:class_cam}
\end{figure}



\begin{table}[t]
\centering
\tabcolsep=0.1cm
\begin{tabular}{@{}lccc@{}}
\toprule
                            & Mixup & Cutout & CutMix \\ \midrule
Usage of full image region  & \textcolor{darkergreen}{\ding{52}}   & \textcolor{red2}{\ding{56}}      & \textcolor{darkergreen}{\ding{52}}      \\
{Regional dropout} & \textcolor{red2}{\ding{56}}     & \textcolor{darkergreen}{\ding{52}}      & \textcolor{darkergreen}{\ding{52}}    \\
Mixed image \& label & \textcolor{darkergreen}{\ding{52}}    & \textcolor{red2}{\ding{56}}     & \textcolor{darkergreen}{\ding{52}}    \\ 
\bottomrule
\end{tabular}
\caption{Comparison among Mixup, Cutout, and CutMix.}
\label{tab:checkbox}
\end{table}



We observe that Cutout successfully lets a model focus on less discriminative parts of the object, such as the belly of Saint Bernard, while being inefficient due to unused pixels.
Mixup, on the other hand, makes full use of pixels, but introduces unnatural artifacts.
The CAM for Mixup, as a result, shows that the model is confused when choosing cues for recognition.
We hypothesize that such confusion leads to its suboptimal performance in classification and localization, as we will see in {Section~\ref{sec:experiments}.}

{CutMix efficiently improves upon Cutout by being able to localize the two object classes accurately. }
We summarize the key differences among Mixup, Cutout, and CutMix in Table~\ref{tab:checkbox}.


\noindent\textbf{Analysis on validation error: }
We analyze the effect of CutMix on stabilizing the training of deep networks.
We compare the top-1 validation error during the training with CutMix against the baseline.
We train ResNet-50~\cite{resnet} for ImageNet Classification, and PyramidNet-200~\cite{pyramidnet} for CIFAR-100 Classification.
Figure~\ref{fig:cls_loss} shows the results.

We observe, first of all, that CutMix achieves lower validation errors than the baseline at the end of training. 
At epoch 150 when the learning rates are reduced, the baselines suffer from overfitting with increasing validation error.
CutMix, on the other hand, shows a steady decrease in validation error; diverse training samples reduce overfitting.